\subsection{Regimes of Leakage}
The averaging law organizes the observations across datasets. In single-trial decoding the leaked feature competes with trial noise, and its effect is small in most paradigms, from 0.4 to 3.7 points. In decoding of averaged data the trial noise shrinks as $1/k$, and the leak approaches the discriminability of the post-stimulus mean itself: 7 to 22 points across participants on ERP CORE and 20 points on SHAPE. Participant-level decoding of averaged ERPs is common in clinical and individual-differences work, where it is also the regime least protected by subject-grouped validation, because the leak lives inside each observation. The law also predicts where the leak is weak. When the condition difference is brief or changes sign within the epoch, as in the MMN, its contribution to the epoch mean is small, and the MMN shows no within-participant leak at any $k$. Across participants its z-scored statistic (0.553, $p=0.045$) sits close to its unscaled value (0.537, $p=0.070$), so its pre-stimulus interval is not a clean null region in the first place.

The prevalence audit places these regimes in current practice. Three of ten widely used codebases scale each sample by its own statistics in their default or reported pipelines, and 70\% of a fixed sample of recent decoding papers leave their input scaling unreported, which leaves the reader without the information needed to assess this leak.
